% Hand-written text and the commit; measured tables come from
% numbers_data.tex, generated by figures/make_numbers.py from the records.
% Generated by figures/make_numbers.py from the records; do not edit.
\newcommand{\SummaryRows}{%
L1 & 1 & no & 1341 & 1344 & 1855 & 5 & 0.20 & 1.20 & 128.4 & 64.5\\
L3 & 3 & no & 1991 & 1992 & 2730 & 84 & 1.14 & 5.83 & 192.6 & 96.5\\
L5 & 5 & yes ($\lambda{=}128$) & 2630 & 2632 & 3594 & 51 & 2.30 & 12.78 & 256.8 & 128.5\\
Q3 & 3 & yes & 3904 & 3904 & 5311 & 815 & 5.88 & 37.34 & 384.4 & 192.5\\
Q5 & 5 & yes & 5192 & 5192 & 7050 & --- & --- & --- & 512.7 & 256.5\\}
\newcommand{\ParamRows}{%
L1 & 128 & 1341 & 21 & 1076 & 162 & 256.8 & 144 & 277 & 1344 & 1855 & 1372\\
L3 & 192 & 1991 & 32 & 1594 & 243 & 385.1 & 1059 & 409 & 1992 & 2730 & 2020\\
L5 & 256 & 2630 & 42 & 2106 & 324 & 513.5 & 380 & 536 & 2632 & 3594 & 2660\\
Q3 & 192 & 3904 & 61 & 3125 & 485 & 768.7 & 371 & 791 & 3904 & 5311 & 3932\\
Q5 & 256 & 5192 & 82 & 4155 & 647 & 1025.5 & 975 & 1049 & 5192 & 7050 & 5220\\}
\newcommand{\PerfRows}{%
L1 (jobs=1) & 4.6 [3.6--6.5], $n{=}5$ & 5.4 [4.5--7.5], $n{=}5$ & 0.204 [0.165--0.304], $n{=}50$ & 1.20 [1.08--1.80], $n{=}50$ & 0.208 / 1.25\\
L3 (jobs=1) & 84.1 [76.5--94.0], $n{=}5$ & 87.8 [79.8--97.3], $n{=}5$ & 1.136 [0.796--1.574], $n{=}25$ & 5.83 [5.40--6.30], $n{=}25$ & 1.038 / 6.24\\
L5 (jobs=1) & 50.6 [39.2--80.3], $n{=}5$ & 60.2 [48.6--89.2], $n{=}5$ & 2.301 [1.744--3.063], $n{=}25$ & 12.78 [11.81--14.90], $n{=}25$ & 2.281 / 13.08\\
Q3 (jobs=4) & 815.3 [689.9--940.6], $n{=}2$ & 2592.3 [2395.0--2789.7], $n{=}2$ & 5.880 [5.381--6.562], $n{=}6$ & 37.34 [36.12--38.49], $n{=}6$ & 5.929 / 37.13\\
Q5 (jobs=?) & --- & --- & --- & --- & --- / ---\\}
\newcommand{\OneCoreRows}{%
L1 & 6.5 [4.6--8.5], $n{=}2$ & 0.240 [0.202--0.341], $n{=}10$ & 1.88 [1.78--2.12], $n{=}10$\\
L3 & 85.2, $n{=}1$ & 1.055 [1.044--1.086], $n{=}3$ & 9.61 [9.33--9.65], $n{=}3$\\
L5 & 135.7, $n{=}1$ & 2.938 [2.879--2.992], $n{=}3$ & 22.20 [21.51--25.18], $n{=}3$\\}
\newcommand{\FieldRows}{%
L1 & 21 & 536 & 1674 & 36,539\\
L3 & 32 & 1760 & 5479 & 59,891\\
L5 & 42 & 2957 & 9019 & 88,919\\
Q3 & 61 & 5936 & 18343 & 152,385\\
Q5 & 82 & 10656 & 32345 & 240,125\\}
\newcommand{\ChainRows}{%
L1 & 569 & 0.63 & 1122 & 670\\
L3 & 828 & 2.99 & 3640 & 664\\
L5 & 1091 & 6.45 & 5922 & 657\\
Q3 & 1598 & 18.69 & 11697 & 638\\}
\newcommand{\BenchRunId}{%
run-20261017T2022Z-kem-bench}
\newcommand{\DecLOne}{%
1.2}
\newcommand{\EncLOne}{%
0.20}
\newcommand{\KgQThree}{%
14\,min}
\newcommand{\StepMults}{%
583--720}
\newcommand{\StepFactor}{%
2.3--2.8}

\newcommand{\EigenxCommit}{d89d790}
\newcommand{\MethodText}{%
One campaign, run \art{\BenchRunId}, measured every operation on one host: a
shared cloud VM with 4 vCPUs (Intel Xeon at 2.10\,GHz, one thread per core),
Linux 6.18, rustc~1.97.0, release profile, \art{eigenx} commit
\art{\EigenxCommit}. Isolation is the lowest level (L0): no core pinning
for the all-core rows, and \texttt{taskset -c 0} for the single-core rows.
Other tenants' load is not controlled. The harness is
\art{eigenx:examples/bench_kem.rs}.
\begin{itemize}[nosep]
\item For each set it draws several fresh keys. Key generation is timed in
  wall time and in process CPU time, because it is a rejection search: a
  single sample says little, and parallel attempts (\texttt{jobs}) trade
  CPU time for wall time.
\item Under each key it runs encapsulation then decapsulation of an honest
  ciphertext several times, one PKE round trip of a 1\,KiB message with
  associated data, and one public-key degree test, the attacker's oracle
  (\S\ref{sec:security}). Every result is checked.
\item In the same binary it times $\F_p$ and $\F_{p^2}$ arithmetic.
\end{itemize}
Samples per set: L1 5 keys $\times$ 10, L3 and L5 5 $\times$ 5, Q3 2
$\times$ 3, Q5 1 $\times$ 3, and single-core L1 2 $\times$ 5, L3 and L5 1
$\times$ 3. The Q sets use \texttt{jobs}=4 for key generation; every other
row uses \texttt{jobs}=1. Medians are reported with the range, because the
samples are few and the host is noisy.}
